\section{Recovering the block integral from the activities}
\label{sec:gaussian-reassembly}
The preceding estimates concern a dictionary of local factors. We now
verify that its absolutely convergent sum is the original block integral.
This verification fixes the scalar normalization and the incidence conditions
before any logarithm of a gas is taken. Throughout, fix the coarse field,
forest, allocation of seeds, and tag array. The funded sector hypotheses,
the chart-access conclusion and the locality statements used in the preceding
sections are retained. In particular, all precisions are positive, the
covariance has its specified finite range, and the same assigned supports
include every variable read by their factors.

\subsection{The Gaussian density and the scalar factor}
Write $A$ for the completed sector precision, $A_0$ for its flat reference,
and $C$ for the localized covariance. All three act on the full real
coordinate space, of dimension $D$. Let $q$ be the colour dimension,
$N$ the number of free forest bonds, and $n$ the number of active free
bonds. The other $N-n$ free bonds are integrated in their original Haar
coordinates. Put
\[
 g=\beta^{-1/2},\qquad b=\sqrt{2\pi},\qquad
 m=c_{\rm H}(bg)^q,
\]
where $c_{\rm H}>0$ is the constant in the local probability-Haar
density of Section~\ref{sec:haar-chart}. With $u=\mu+gz$, the reference
Gaussian has density
\begin{equation}\label{eq:reassembly-gaussian-density}
 \vartheta_C(z)=b^{-D}
       \exp\left(-\tfrac12\log\det C-\tfrac12 z^TC^{-1}z\right).
\end{equation}
Inactive coordinates are completed by independent standard Gaussians.
Their energy is $\mathscr G(z)=\frac12\sum_{i\ {\rm inactive}}z_i^2$.
There are $D-nq$ such coordinates, including the tree coordinates.
Consequently the change of variables in the active Haar coordinates,
with the dummy Gaussian integrals inserted, gives the multiplier
\begin{equation}\label{eq:reassembly-jacobian}
 b^{-(D-nq)}g^{nq}c_{\rm H}^{n}
       \exp\left(-\mathscr G(z)+\sum_{e\ {\rm active}}j_e(u)\right).
\end{equation}
Here $j_e$ is the localized logarithmic Haar correction; its auxiliary
cutoff agrees with the actual correction wherever the smaller chart mask
is nonzero. The constant in \eqref{eq:reassembly-jacobian} is exactly
\begin{equation}\label{eq:reassembly-count}
 b^{-D}m^N\exp(-(N-n)\log m).
\end{equation}
Indeed the right side is $b^{-D}m^n=b^{-(D-nq)}g^{nq}c_{\rm H}^n$.
This uses $n+(N-n)=N$, not the number of active coordinates as a
substitute for the number of free bonds.

The dummy insertion is exact even though the active integrand depends
on the remaining Haar coordinates. Split the latter product measure
into active and inactive free bonds and freeze the inactive variables.
The active affine substitution has Jacobian $g^{nq}$. In the orthogonal
coordinate decomposition into active and dummy real variables, the
dummy density integrates to one. Fubini therefore proves the identity
for each fixed inactive configuration and then for its Haar integral.
Chart access guarantees that the nonzero original integrand is covered
by this substitution; the smaller mask can be inserted there with value
one. No assertion that the entire Haar group lies in one chart is used.

Let $\mathscr L$ denote the original positive prefactor: it is the
noise normalization times $\exp(v|\Lambda_f|)$ times the coarse-field
scalar extracted from the inherited exponent. The scalar retained
outside all activities is
\begin{equation}\label{eq:reassembly-overhead}
                 \mathscr O=\mathscr L\,m^N(\det A_0)^{-1/2}>0.
\end{equation}
The flat reference uses the very same forest. Its determinant is not
replaced by a componentwise determinant, nor is \eqref{eq:reassembly-overhead}
charged to every component.

\subsection{Restoring the masks and the exponent}
We spell out the cancellations that identify the density. Let
$\mathscr E(u)$ be the whole quadratic-plus-remainder block energy,
and let $r_j(u)$ be the inherited logarithms. Let $d_p(u)$ be the
individual nonlinear remainders, so the signed channel average satisfies
\[
          \mathscr E(u)=\langle\textstyle\sum_p d_p(u)\rangle
                                      +Q(u),
\]
where $Q(u)=\frac12\|Mu+b_{\rm row}\|^2$ is the physical quadratic
expression. The selected row matrix $M$ has zero inactive columns;
$A=M^TM+I_{\rm inactive}$ and $\ell=M^Tb_{\rm row}$. The inactive
coordinates of $\mu$ and $\ell$ are zero. This is the rowwise square expansion used in
Sections~\ref{sec:mean-action-jets} and \ref{sec:remote-action-jets}:
each row is assigned either to its seed neighbourhood or to the remote
remainder, with no omitted row. Write $c_e=\frac12(J_e(A_0)-J_e(A))$
for the local normalization correction.

The compulsory packet exponent, summed over the packet layout, is
\begin{equation}\label{eq:reassembly-packet-log}
 s(u)=-\beta\bigl(D_{\rm near}(u)+Q_{\rm seed}(\mu)\bigr)
       +\sum_{e\ {\rm near}}c_e+\sum_{e\ {\rm near}}j_e(u)
       -(N-n)\log m.
\end{equation}
Here $D_{\rm near}(u)$ is the signed average of $d_p(u)$ over the
posted rows at radius $200B$, and $Q_{\rm seed}(\mu)$ is the signed
average of the extracted row squares at $\mu$ over those same posted
rows. They are not complementary row sets. The two occurrences of
``near'' in the following local sums use their respective radii, $200B$
and $304B$; they are different sets. The packet layout uses disjoint
halos of radius at least $500B$, so all reads through $304B$, including
the extra two lattice steps, belong to a unique packet. Thus summing
the packet exponents gives exactly \eqref{eq:reassembly-packet-log}.
Likewise the packet products reproduce the compulsory likelihoods,
flag tests and chart masks once each. A test whose location belongs to
a packet has the same seed membership in that packet as in the full
allocation: a seed changing the test would meet that packet's halo,
and hence would lie in the same component.

Let $f_i(z)$ be the optional logarithms in the complete dictionary.
On the nonzero smaller chart mask their sum is
\begin{align}
 \sum_i f_i(z)={}&\sum_j r_j(u)+\sum_p h_p(u)
                +\sum_{e\ {\rm remote}}c_e
                +\sum_{e\ {\rm remote}}j_e(u)+\mathscr R(z),
                                                        \label{eq:reassembly-logs}\\
 \mathscr R(z)={}&\tfrac12(\log\det C+\log\det A_0)
           -\sum_e c_e+z^T\widehat h
                       +\tfrac12z^T(C^{-1}-A)z.\nonumber
\end{align}
The remote index sets here are the exact complements of the two near
sets above. The determinant and inverse series giving $\mathscr R$
are those of Sections~\ref{sec:local-covariance} and
\ref{sec:itinerary-prices}. Their three determinant contributions are
\[
 -\tfrac12(\log\det A-\sum_eJ_e(A)),\quad
 \tfrac12(\log\det A_0-\sum_eJ_e(A_0)),\quad
 \tfrac12(\log\det A+\log\det C).
\]
Their sum is the determinant part of $\mathscr R$. The linear
itineraries sum to $z^T\widehat h$ and the quadratic itineraries sum
to $\frac12z^T(C^{-1}-A)z$. Their endpoint protectors equal one on
the smaller full mask. Absolute summability was proved before this
recombination, so these are identities of convergent series.

The remote action terms obey
\begin{equation}\label{eq:reassembly-remote-cancel}
 \sum_p h_p(u)-\beta\bigl(D_{\rm near}(u)+Q_{\rm seed}(\mu)\bigr)
       =-\beta\bigl(\langle\textstyle\sum_pd_p(u)\rangle+Q(\mu)\bigr).
\end{equation}
For a near row this follows from its assigned debit and mean-square
term. For a remote row it is precisely the unmasked local-square
identity: its protecting cutoff equals one on the smaller full mask.
Summing first over rows and then over the finite signed channels proves
\eqref{eq:reassembly-remote-cancel}, including rows with zero factors.

Finally, if $\ell$ is the linear coefficient of $Q$, then
$\widehat h=-g^{-1}(A\mu+\ell)$ and $\beta g^2=1$ give
\begin{equation}\label{eq:reassembly-square-cancel}
 -\beta Q(u)-\mathscr G(z)
 =-\beta Q(\mu)+z^T\widehat h
     +\tfrac12z^T(C^{-1}-A)z-\tfrac12z^TC^{-1}z.
\end{equation}
To verify this, expand $Q(\mu+gz)$, use that the inactive mean and
linear coefficient vanish, and add the inactive identity quadratic
form supplied by $\mathscr G$. This is an algebraic identity, valid
also where the masks vanish.

Combining \eqref{eq:reassembly-packet-log}--\eqref{eq:reassembly-square-cancel}
cancels both local sums $\sum c_e$ and both full determinants. It gives
\begin{align}
 &-\tfrac12\log\det A_0+s(u)+\sum_i f_i(z)
             -\tfrac12\log\det C-\tfrac12z^TC^{-1}z\nonumber\\
 &\qquad=\sum_jr_j(u)-\beta\mathscr E(u)-\mathscr G(z)
          +\sum_{e\ {\rm active}}j_e(u)-(N-n)\log m.
                                                    \label{eq:reassembly-exponents}
\end{align}
Multiplication by $b^{-D}m^N\mathscr L$ and
\eqref{eq:reassembly-count} now recover the exact original density
and its Haar Jacobian.

There is also a finite product identity for the masks. Write $a_p$
for a compulsory test and $w_p$ for its optional warning. At a
compulsory position $w_p=0$ and its packet supplies $a_p$; at every
other position the optional factor is $1+w_p=a_p$ on the smaller
full mask. The auxiliary mask warnings restore that full mask from
the packet masks. If it vanishes, the entire product on both sides
is zero. If it is nonzero, every doubled mask is one and the preceding
identities apply literally. The complete product of physical tag
likelihoods equals the number of tag choices times the original
normalized likelihood. Thus a single factor $|\mathcal T|^{-1}$,
where $\mathcal T$ is the tag-array space, remains outside the product.
This is the uniform tag average, not a new likelihood for each use
of a tag.

\subsection{Absolute expansion and return to probability Haar}
Fix a finite old forced-label set $K$ and a packet layout $\mathcal P$.
For finite sets $I$ of logarithmic labels and $J$ of warning labels,
the expanded integrand is
\begin{equation}\label{eq:reassembly-molecule}
 \Phi_{K,I,J}
  =\prod_{P\in\mathcal P}P(u)\prod_{k\in K}K_k(u)
       \prod_{i\in I}(e^{f_i(z)}-1)\prod_{j\in J}w_j(u).
\end{equation}
The four alphabets form a disjoint union; coincident physical locations
do not identify their labels. On the finite lattice the nonempty
supports and rooted price bounds imply $\sum_i p_i<\infty$ for the
whole alphabet. Separately, the all-argument logarithm estimates of
Sections~\ref{sec:remote-action-jets}--\ref{sec:inherited-activity-prices}
give envelopes $|f_i(z)|\le a_i$ with $\sum_i a_i<\infty$. Their
endpoint cutoffs are retained in this assertion. Thus
$\sum_i|f_i(z)|<\infty$, and
\[
 \exp\left(\sum_i f_i(z)\right)
                  =\sum_{I\ {\rm finite}}\prod_{i\in I}(e^{f_i(z)}-1).
\]
For a proof, exhaust the countable alphabet by finite subsets and
expand each finite product. Absolute convergence follows from
$\sum_i|e^{f_i}-1|<\infty$: outside finitely many indices,
$|e^{f_i}-1|\le2|f_i|$. The warning alphabet is finite and its
expansion is the ordinary identity for $\prod_j(1+w_j)$.

The Gaussian moment estimate gives, with the original probability
Haar variables $y$ retained,
\[
 \int|\Phi_{K,I,J}(y,z)|\,d\nu_F(y)\,d\gamma_C(z)
                         \le\prod_{v\in\mathcal P\sqcup K\sqcup I\sqcup J}p_v.
\]
Summing the right side over all finite label sets gives at most
$\prod_v(1+p_v)\le\exp(\sum_vp_v)$. Distinct triples $(K,I,J)$
give distinct label sets. Tonelli applied to the absolute values
therefore justifies every exchange between these sums and integrals.
The same argument covers the subsequent finite sum over seed
allocations and tag arrays. There is no division by the value of a
vacuum or flagged term.

An allocation contributes only when its flags match the coarse field,
the old forced labels are admissible, and their inventories have the
required incidence with the seeds. These conditions are necessary for
a nonzero original integrand: the profile enforces its exact flags,
and a nonzero old gas term enforces its admissibility and incidence.
Set the entire integral to zero when any condition fails. On a valid
allocation the density calculation above applies; on an invalid one
both sides vanish. Hence the absolutely convergent sum, multiplied by
$|\mathcal T|^{-1}\mathscr O$, equals the original contribution for
the fixed tags and seed allocation. Summing the old forced sets,
then the tags and seed allocations, recovers the probability-Haar
block push with scalar factor \eqref{eq:reassembly-overhead}.

\subsection{Factorization into connected activities}
It remains to show that the integral of a molecule factors over
its connected components. Connectivity is defined by intersection
of the full assigned supports. For fixed tags, different components
read disjoint Haar coordinates. Their Gaussian coordinates may be
larger than their physical read sets; the assigned supports include
the prescribed covariance buffer. If two components have disjoint
assigned supports, any two Gaussian coordinates read by them are
farther apart than the range of $C$. Their cross covariance is zero.
The Gaussian characteristic function consequently splits into the
product of its component characteristic functions, proving independence
of these coordinate projections. It follows that the Gaussian integral
of the product factors. The remaining product-Haar integral factors
by its disjoint read sets. All integrands are measurable and integrable
by the preceding absolute bounds.

Lastly average over the tags. A component integral depends only on
the site tags and on the bond tags anchored in its assigned support.
These coordinate sets are disjoint for different components. The
uniform measure on tag arrays is a product measure, so averaging first
the bond tags and then the site tags proves the same factorization
for the averaged integrals. Averages are taken after each component
has been identified, using the common original tag space.

For each connected label set $M$, define its activity $a_M(V)$ by
this averaged integral with the exact local incidence tests: its
inventory must equal the flags of $V$ inside its assigned support,
and all its forced-label and packet tests must hold. The integral
uses the globally flagged sector with the same local data. Locality
shows that its value agrees with the restriction of the full
allocation: on this support the seeds, tags, forced labels and coarse
field coincide. The flags outside the support cannot change its value.

Taking connected components and taking their union are inverse on
admissible label families. Indeed every label has nonempty support,
so distinct components cannot share a label. Their supports are
disjoint by definition. Exact global flag coverage is equivalent to
exact local inventory equality in each component together with
coverage by all inventories. This proves the weight-preserving
factorization, including the empty family. Reindexing the absolutely
convergent sum gives
\begin{equation}\label{eq:reassembly-block-gas}
 (\mathcal B R)(V)=\mathscr O(V)
   \sum_{\substack{\mathcal A\ {\rm finite,\ compatible}\\
                     \bigcup_{M\in\mathcal A}I_M=\operatorname{Flags}(V)}}
                         \prod_{M\in\mathcal A}a_M(V).
\end{equation}
Thus the local prices bound an exact representation of the original
block push. The later separation of flag-free logarithms and the
iteration of the resulting bounds still use their stated construction
inputs; neither follows merely from \eqref{eq:reassembly-block-gas}.
